% =====================================================================
%  1.4 Rasjonale uttrykk — fasit
% =====================================================================

\fk{Del 1 — Forkorte, gange og dele}
\fo{1.1} a) $\frac{2x}{3}$, $x \neq 0$ \ b) $3$, $x \neq -2$ \ c) $x-5$, $x \neq 0$ \ d) $\frac{1}{x}$, $x \neq 0$ og $x \neq 4$
\fo{1.2} a) $x^2-25 = (x-5)(x+5)$ og $x^2-5x = x(x-5)$ \ b) $\frac{x+5}{x}$, $x \neq 0$ og $x \neq 5$
\fo{1.3} a) $x \neq 3$ \ b) $x \neq 2$ og $x \neq -2$ \ c) $x \neq 0$ og $x \neq -2$
\fo{1.4} a) $\frac{x-3}{x+3}$, $x \neq \pm 3$ \ b) $\frac{2(x-2)}{x+2}$, $x \neq -2$ \ c) $\frac{x+3}{3}$, $x \neq -3$ \ d) $-1$, $x \neq 5$
\fo{1.5} a) $2(x-1)$, $x \neq 0$ og $x \neq -1$
b) $\frac{(x-3)(x+3)}{x+1}\cdot\frac{2(x+1)}{x+3} = 2(x-3)$, $x \neq -1$ og $x \neq -3$
c) $\frac{x(x+2)}{(x-2)(x+2)}\cdot\frac{x-2}{x^2} = \frac{1}{x}$, $x \neq 0$, $x \neq 2$ og $x \neq -2$
\fo{1.6} Feilen er i første pil: $x^2$ og $-9$ er ledd, og ledd kan ikke forkortes. Riktig:
$\frac{(x-3)(x+3)}{x(x-3)} = \frac{x+3}{x}$, $x \neq 0$ og $x \neq 3$.
\fo{1.7} a) Noen ganger. Stemmer for alle $x \neq 1$, men venstre side er ikke definert for $x = 1$.
b) Alltid (når $x \neq 3$), fordi $3-x = -(x-3)$.
c) Noen ganger. $5(x+3) = 3(x+5)$ gir $2x = 0$, så den stemmer bare for $x = 0$.
\fo{1.8} Gang teller og nevner med $x-3$: $\frac{(x+1)(x-3)}{(x-2)(x-3)}$. Forkortet er det $\frac{x+1}{x-2}$.
Vilkår: $x \neq 2$ og $x \neq 3$. Flere svar er riktige.
\fo{1.9} Det opprinnelige uttrykket har $x-2$ i nevneren, så det er ikke definert for $x = 2$.
Vilkåret kommer fra det opprinnelige uttrykket og gjelder fortsatt etter forkorting.
Det forkortede uttrykket gir $\frac{2}{4}$ for $x = 2$, men det er et annet uttrykk.
\fo{Sjekk 1} a) $\frac{2}{x-2}$, $x \neq \pm 2$ \ b) $x \neq 5$ \ c) For $x = 1$ er $\frac{4}{6} = \frac{2}{3}$, men $\frac{3}{5}$ er noe annet. Tre og fem er ledd, ikke faktorer.

\fk{Del 2 — Addere og subtrahere}
\fo{2.1} a) $\frac{7}{12}$ \ b) $\frac{11}{24}$ \ c) $\frac{7}{10}$
\fo{2.2} a) $\frac{8}{x}$ \ b) $\frac{3x+6}{x+3} = \frac{3(x+2)}{x+3}$ \ c) $\frac{3}{x}$ \ d) $\frac{3}{x-1}$
\fo{2.3} a) $x(x+2)$ \ b) $x^2-4 = (x-2)(x+2)$ \ c) $6x^2$
\fo{2.4} a) $\frac{3x+3}{x(x+3)} = \frac{3(x+1)}{x(x+3)}$ \ b) $\frac{x+2}{x(x-1)}$ \ c) $\frac{3x+2}{6x^2}$ \ d) $\frac{x^2+1}{(x+1)(x-1)}$
\fo{2.5} a) $\frac{x+3-6}{(x-3)(x+3)} = \frac{1}{x+3}$, $x \neq \pm 3$
b) $\frac{2x+(x-1)}{(x-1)(x+1)} = \frac{3x-1}{(x-1)(x+1)}$, $x \neq \pm 1$
c) $\frac{(x+1)-1}{x(x+1)} = \frac{1}{x+1}$, $x \neq 0$ og $x \neq -1$
\fo{2.6} Feilen er i første pil: minuset gjelder hele $x-1$. Riktig: $\frac{2-(x-1)}{x} = \frac{2-x+1}{x} = \frac{3-x}{x}$, $x \neq 0$.
\fo{2.7} a) Tid er strekning delt på fart. Summen er $\frac{2s+3s}{120} = \frac{5s}{120} = \frac{s}{24}$ timer.
b) Total strekning er $2s$. Farten er $\frac{2s}{s/24} = 48$ km/h, ikke 50 km/h. Hun bruker mer tid på den langsomme turen.
\fo{2.8} a) Ingen løsning. Ganger vi med $x-2$, får vi $x = 2+3(x-2)$, så $x = 2$. Men $x = 2$ er ikke tillatt.
b) Stemmer for alle $x \neq 1$. Venstre side er $\frac{2+(x-1)}{x-1} = \frac{x+1}{x-1}$, som er høyre side.
\fo{2.9} Venstre side er $\frac{a+b}{ab}$ og høyre side er $\frac{2}{a+b}$. Kryssmultiplisering gir $(a+b)^2 = 2ab$, altså $a^2+b^2 = 0$.
En sum av kvadrater er null bare når $a = b = 0$, og det er ikke tillatt. Svar: aldri.
\fo{Sjekk 2} a) $\frac{x+y}{xy}$, $x \neq 0$, $y \neq 0$ \ b) $\frac{2-x}{x}$, $x \neq 0$ \ c) Nevnerne skal ikke legges sammen. $\frac{1}{2}+\frac{1}{3} = \frac{3}{6}+\frac{2}{6} = \frac{5}{6}$.

\fk{Ekstra utfordring}
\fo{U1} $\frac{(x-2)(x+2)}{(x+2)^2}\cdot\frac{x(x+2)}{x-2} = x$. Vilkår: $x \neq -2$, $x \neq 0$ og $x \neq 2$.
\fo{U2} Felles nevner $(x-1)(x+1)$: $\frac{(x+1)-(x-1)-2}{x^2-1} = \frac{0}{x^2-1} = 0$. Vilkår: $x \neq \pm 1$.
Uttrykket er lik 0 for alle tillatte $x$.
\fo{U3} Teller: $\frac{x+1}{x}$. Nevner: $\frac{x^2-1}{x^2}$. Da er $\frac{x+1}{x}\cdot\frac{x^2}{(x-1)(x+1)} = \frac{x}{x-1}$. Vilkår: $x \neq 0$, $x \neq \pm 1$.
\fo{U4} a) $x+1+\frac{2}{x-1} = \frac{(x+1)(x-1)+2}{x-1} = \frac{x^2-1+2}{x-1} = \frac{x^2+1}{x-1}$.
b) Leddet $\frac{2}{x-1}$ blir nær 0 når $x$ er stor. Uttrykket er da omtrent $x+1$.
\fo{U5} a) $\frac{1}{n}-\frac{1}{n+1} = \frac{n+1-n}{n(n+1)} = \frac{1}{n(n+1)}$.
b) Hvert ledd er $\frac{1}{n}-\frac{1}{n+1}$. Summen blir $\left(1-\frac12\right)+\left(\frac12-\frac13\right)+\ldots+\left(\frac{1}{99}-\frac{1}{100}\right)$.
Alt midt i forsvinner, og det er $1-\frac{1}{100} = \frac{99}{100}$.
\fo{U6} a) Summen for $n$ ledd er $\frac{n}{n+1}$: $\frac12, \frac23, \frac34, \ldots, \frac{10}{11}$.
b) Programmet viser bare ti tilfeller. Det sier ikke at mønsteret gjelder for alle $n$. Regningen i U5 a) gjelder for alle $n$, og er et bevis.
\fo{U7} $\frac{x+5}{x-1} = 1+\frac{6}{x-1}$. Brøken er et helt tall når $x-1$ går opp i 6, altså $x-1 \in \{\pm1, \pm2, \pm3, \pm6\}$.
Da er $x \in \{-5, -2, -1, 0, 2, 3, 4, 7\}$. Vi må ha $x \neq 1$.
